\section{Open-World Execution and Recovery}
\label{sec:rw_execution}

Even when perception and planning succeed, action execution in open-world environments may fail due to unexpected dynamics, novel object properties, or environmental changes. Recovery from such failures requires mechanisms that go beyond standard error handling: the agent must identify why the failure occurred and adapt its behavior accordingly.

\paragraph{Plan monitoring and repair.}
Classical approaches to execution-time adaptation include plan monitoring and repair \citep{fox2006plan, fritz2007monitoring}, which detect deviations from expected execution outcomes and revise plans within the existing domain model. However, plan repair assumes that the domain model remains valid and that a correct plan exists within the current model, which may not hold when qualitative novelty has occurred. The cascading recovery pipeline in \cref{ch:rapid_learn} extends this idea: it attempts lightweight symbolic recovery first (analogous to plan repair) and escalates to learning-based executor discovery only when symbolic methods are insufficient. This cascade reflects the insight that not all failures require new learning---some can be resolved through replanning within the existing model.

\paragraph{Sim-to-real transfer and domain adaptation.}
Sim-to-real transfer techniques, including domain randomization \citep{tobin2017domain} and system identification \citep{peng2018sim}, improve policy robustness by training across varied conditions, but the variations are typically parameterized in advance and may not cover the kinds of qualitative changes encountered in open-world settings. The interaction abstractions developed in \cref{ch:generalization} adopt a complementary strategy: rather than randomizing over the robot's parameters to achieve robustness, they remove the robot from the learning problem entirely by operating in object-centric force space. This means that sim-to-real transfer requires only a mapping from learned forces to robot commands, rather than bridging a full embodiment gap.

\paragraph{Continual and lifelong learning.}
Continual learning methods \citep{kirkpatrick2017overcoming, rusu2016progressive} address the problem of acquiring new capabilities without forgetting previous ones, but most assume known task boundaries and distributions. In the open-world setting, the agent does not know in advance when novelty will occur or what form it will take. The novelty detection and characterization mechanisms studied in \cref{ch:benchmarks} and \cref{ch:rapid_learn} address this gap by providing the agent with the ability to detect mismatch between expected and observed execution outcomes, triggering adaptation only when necessary.

\paragraph{Integrated approaches.}
Some efforts have sought to address multiple stages of the open-world agent loop in an integrated manner. World models \citep{ha2018world, hafner2023mastering} offer a potential unifying framework by learning environment dynamics that can serve both perception and planning, though they typically do not handle discrete, qualitative changes in the environment. Task and motion planning systems \citep{garrett2021integrated, kaelbling2011hierarchical} integrate symbolic and geometric reasoning but inherit the brittle model assumptions of classical planning. Despite growing interest in open-world autonomy, the field remains largely fragmented: methods that address perception under novelty, planning under incomplete models, and recovery from execution failures are typically developed in isolation from one another. The framework developed in this dissertation takes a step toward integration by connecting symbolic planning (for task decomposition and failure localization) with reinforcement learning (for behavior acquisition) through structured abstractions that mediate between the two levels.